% =====================================================================
%  FASE DE EXPLORACIÓN (versión actualizada). Cambios respecto al original:
%   - Realidad aumentada: se corrige «el usuario siempre confirma la máquina»
%     (con confianza >= 0,8 la app acepta la sugerencia) y se acota el término
%     según Azuma (1997).
%   - Actores y tecnologías: Cloud Functions y App Check; USDA vía función.
%   - Fuentes de datos: se añaden ExerciseDB y Gym Visual; «s.f.» -> «s. f.».
%   - RNF1: mínimo de Android = 7.0 (API 24), valor por defecto de React Native.
%   - Persistencia sin conexión: capacidad del SDK, no activada en la app.
%   - Cronograma: solicitudes de integración #5, #6 y #7.
% =====================================================================
% =====================================================================
\section{Fase de Exploración}
\label{sec:des-exploracion}

En esta fase se definen los requerimientos, el alcance y las tecnologías para el desarrollo
del producto.

\subsection{Establecimiento de los interesados}

Tutor del proyecto(Cliente): Ing. Jimmy Villarroel Novillo.

\subsection{Actores del sistema}

El actor principal del sistema es el usuario, quien completa su perfil,
consulta las instrucciones para el uso correcto de las máquinas del gimnasio y genera 
rutinas de entrenamiento personalizadas. Además, puede registrar su consumo
calórico, consultar su índice de masa corporal (IMC) y hacer un seguimiento
de su progreso físico.

\subsection{Roles y Tareas}

Para el desarrollo del sistema se han identificado los siguientes usuarios:

\begin{table}[htbp]
\centering
\begin{tabular}{|p{0.18\linewidth}|p{0.49\linewidth}|p{0.21\linewidth}|}
  \hline
  \textbf{Nombre} & \textbf{Descripción} & \textbf{Stakeholder} \\
  \hline
  Usuario Final & 
  Completa su perfil físico, consulta las instrucciones para el uso correcto
  máquinas del gimnasio, escanea máquinas, obtiene y planifica sus rutina, registra 
  su consumo calórico, consulta su índice de masa corporal IMC y revisa su progreso físico. & 
  Persona que entrena en el gimnasio\\
  \hline
  \end{tabular}
  \caption{Roles y tareas de Coach App (Elaboración propia)}
  \label{tab:roles-tareas}
\end{table}


\subsection{Requerimientos iniciales} 

En esta etapa se elabora la primera definición de los requerimientos. Los requerimientos 
funcionales describen las acciones que la aplicación debe permitir al usuario, mientras que
los no funcionales establecen las condiciones de calidad que debe cumplir.
Esta definición inicial delimita el alcance del sistema y sirve de base para su diseño y desarrollo.

\subsubsection{Requerimientos funcionales}

Los requerimientos funcionales describen las funciones que la aplicación debe ofrecer al usuario. 
La siguiente tabla presenta los requerimientos generales y sus respectivos subrequerimientos.

\begin{longtable}{|L{2.2cm}|L{10.8cm}|}
\hline
\rowcolor{encabezado}
\textbf{\textcolor{black}{CÓDIGO}} &
\textbf{\textcolor{black}{DESCRIPCIÓN}} \\
\hline
\endfirsthead

\hline
\rowcolor{encabezado}
\textbf{\textcolor{black}{CÓDIGO}} &
\textbf{\textcolor{black}{DESCRIPCIÓN}} \\
\hline
\endhead

\endfoot

\caption{Requerimientos funcionales del sistema (elaboración propia)}
\label{tab:rf}\\
\endlastfoot

RF1.1 & Permitir el registro de un usuario con nombre, apellido,
correo electrónico y contraseña. \\
\hline
RF1.2 & Permitir el inicio de sesión mediante correo electrónico y
contraseña o mediante una cuenta de Google. \\
\hline
RF1.3 & Permitir la recuperación de la contraseña por correo
electrónico y el cierre de sesión. \\
\hline
RF1.4 & Permitir que el usuario registre, en su primer ingreso, su
objetivo de entrenamiento, nivel de experiencia, edad, estatura, peso
y días disponibles para entrenar. \\
\hline
RF1.5 & Permitir la consulta y edición de los datos del perfil. \\
\hline

RF2.1 & Permitir la búsqueda de máquinas por texto, con o sin tildes,
y el filtrado del catálogo por grupo muscular y tipo de equipo. \\
\hline
RF2.2 & Mostrar la ficha de cada máquina con su descripción, los
músculos trabajados, las instrucciones de uso, los errores comunes,
las recomendaciones para prevenir lesiones y las variantes comerciales. \\
\hline
RF2.3 & Mostrar los ejercicios asociados a cada máquina con su
animación, músculo objetivo y nivel de dificultad. \\
\hline
RF2.4 & Permitir el escaneo de una máquina con la cámara y mostrar
su ficha sobre la imagen en vivo. Si el sistema no logra identificarla
automáticamente, permitir que el usuario la seleccione de una lista. \\
\hline

RF3.1 & Generar una rutina semanal según el objetivo de entrenamiento
y el nivel de experiencia del usuario. \\
\hline
RF3.2 & Permitir la creación y eliminación de rutinas personalizadas
organizadas por días, con series, repeticiones y tiempos de descanso
para cada ejercicio. \\
\hline
RF3.3 & Permitir que el usuario agregue ejercicios a su rutina desde
la ficha de una máquina y marque una rutina como activa. \\
\hline

RF4.1 & Calcular el índice de masa corporal (IMC), clasificar el
resultado según los criterios de la OMS y guardar el historial de
cálculos. \\
\hline
RF4.2 & Permitir el registro del peso, la estatura y las notas de
progreso físico, con cálculo automático del IMC. \\
\hline

RF5.1 & Permitir la búsqueda de alimentos en español y el registro
de la cantidad consumida en gramos. Calcular las calorías y los
macronutrientes de la porción registrada. \\
\hline
RF5.2 & Permitir que el usuario agregue alimentos de consumo frecuente
con un toque o registre un alimento manualmente. \\
\hline
RF5.3 & Mostrar el consumo diario de calorías y macronutrientes en
relación con las metas del usuario y permitir la eliminación de
registros. \\
\hline

RF6.1 & Mostrar en un gráfico la evolución del peso y del IMC, e
indicar la variación acumulada y la categoría actual del IMC. \\
\hline
RF6.2 & Permitir que el usuario marque como completado el entrenamiento
del día y consulte su racha, los días entrenados durante la semana y
el total de entrenamientos completados. \\
\hline
\end{longtable}

\subsubsection{Requerimientos no funcionales}

Los requerimientos no funcionales establecen las condiciones de calidad
que debe cumplir la aplicación. La tabla siguiente describe cada requerimiento.

\begin{longtable}{|L{2.2cm}|L{10.8cm}|}
\hline
\rowcolor{encabezado}
\textbf{\textcolor{black}{CÓDIGO}} &
\textbf{\textcolor{black}{DESCRIPCIÓN}} \\
\hline
\endfirsthead

\hline
\rowcolor{encabezado}
\textbf{\textcolor{black}{CÓDIGO}} &
\textbf{\textcolor{black}{DESCRIPCIÓN}} \\
\hline
\endhead

\endfoot

\caption{Requerimientos no funcionales del sistema (elaboración propia)}
\label{tab:rnf}\\
\endlastfoot

RNF1 & \textbf{Portabilidad.} La aplicación debe poder instalarse y
ejecutarse en dispositivos con Android 7.0 (API 24) o superior. \\
\hline

RNF2 & \textbf{Usabilidad.} La interfaz debe presentarse en español,
adaptarse a distintos tamaños de pantalla y permitir el acceso a los
módulos principales mediante pestañas de navegación. \\
\hline

RNF3 & \textbf{Fiabilidad.} Los cambios guardados en las rutinas, los
registros y los entrenamientos deben sincronizarse mediante Firebase
y reflejarse en pantalla sin que el usuario tenga que recargarla
manualmente. \\
\hline

RNF4 & \textbf{Seguridad.} Cada usuario debe poder consultar y modificar
únicamente sus propios datos. El catálogo debe ser de solo lectura
para los usuarios de la aplicación. Las reglas de Firestore deben
validar los tipos y rangos de los datos al guardar los cambios. \\
\hline

RNF5 & \textbf{Eficiencia.} El catálogo de máquinas debe descargarse
una sola vez por sesión y reutilizarse al volver a consultarlo.
Las URL de las imágenes deben conservarse en caché para reducir
las solicitudes repetidas. \\
\hline

RNF6 & \textbf{Mantenibilidad.} El código debe escribirse en TypeScript
con el modo estricto habilitado y organizarse según las capas del
patrón MVC. La comprobación de tipos y el análisis estático deben
completarse sin errores. \\
\hline
\end{longtable}

\subsubsection{Análisis de los requerimientos mediante casos de uso}

La Figura~\ref{fig:casos-uso} presenta el diagrama de casos de uso del
sistema. El actor principal es el usuario; Firebase y las APIs externas
participan como actores secundarios.

\begin{figure}[!t]
\centering
\begin{tikzpicture}[font=\footnotesize, lbl/.style={font=\scriptsize, fill=white, inner sep=1pt}]
  \actor{usr}{(0,-0.4)}{Usuario}
  \actor{fb}{(13.8,2.2)}{Firebase}
  \actor{api}{(13.8,-3.0)}{APIs externas}
  \node[draw, rounded corners, minimum width=9.8cm, minimum height=13.6cm,
        label={[font=\small\bfseries]above:Sistema \nombreApp}] (sis) at (6.75,-0.4) {};
  \node[casouso] (u1)  at (4.2, 5.6) {Registrarse / iniciar sesión};
  \node[casouso] (u2)  at (4.2, 4.4) {Completar \emph{onboarding}};
  \node[casouso] (u3)  at (4.2, 3.2) {Consultar guía de máquinas};
  \node[casouso] (u4)  at (4.2, 2.0) {Escanear máquina (RA)};
  \node[casouso] (u5)  at (4.2, 0.8) {Generar rutina automática};
  \node[casouso] (u6)  at (4.2,-0.4) {Crear rutina personalizada};
  \node[casouso] (u7)  at (4.2,-1.6) {Calcular IMC};
  \node[casouso] (u8)  at (4.2,-2.8) {Registrar progreso físico};
  \node[casouso] (u9)  at (4.2,-4.0) {Registrar alimentos};
  \node[casouso] (u10) at (4.2,-5.2) {Marcar entrenamiento};
  \node[casouso] (u11) at (4.2,-6.4) {Gestionar perfil};
  \node[casouso] (v1) at (9.3, 3.2) {Ver ficha y ejercicios};
  \node[casouso] (v2) at (9.3, 2.0) {Agregar ejercicio a rutina};
  \node[casouso] (v7) at (9.3, 0.8) {Obtener ejercicios (wger)};
  \node[casouso] (v3) at (9.3,-0.4) {Activar rutina};
  \node[casouso] (v5) at (9.3,-2.8) {Ver reporte de progreso};
  \node[casouso] (v4) at (9.3,-4.0) {Buscar alimento (USDA)};
  \node[casouso] (v6) at (9.3,-5.2) {Ver racha semanal};
  \foreach \u in {u1,u2,u3,u4,u5,u6,u7,u8,u9,u10,u11} \draw (usr) -- (\u.west);
  \draw[flechad] (u3) -- node[lbl, above=1pt]{«include»} (v1);
  \draw[flechad] (v2) -- node[lbl, right=1pt]{«extend»} (v1);
  \draw[flechad] (u5) -- node[lbl, above=1pt]{«include»} (v7);
  \draw[flechad] (v3) -- node[lbl, above=1pt]{«extend»} (u6);
  \draw[flechad] (v5) -- node[lbl, above=1pt]{«extend»} (u8);
  \draw[flechad] (u9) -- node[lbl, above=1pt]{«include»} (v4);
  \draw[flechad] (v6) -- node[lbl, above=1pt]{«extend»} (u10);
  \draw (v1.east) -- (fb);
  \draw (v5.east) -- (fb);
  \draw (v7.east) -- (api);
  \draw (v4.east) -- (api);
\end{tikzpicture}
\caption{Diagrama de casos de uso del sistema}
\label{fig:casos-uso}
\fuente{Elaboración propia.}
\end{figure}

\clearpage
\subsection{Definición del alcance}
\label{sec:des-alcance}

El alcance del proyecto se determina a través de sus limitaciones:

\subsubsection{Limitaciones}

Las limitaciones de la aplicación desarrollada son:

\begin{itemize}
    \item La aplicación solo puede ser ejecutada en dispositivos con el sistema operativo Android.

    \item El reconocimiento de máquinas en la nube y la búsqueda de
    alimentos requieren conexión a internet.

    \item El escaneo de máquinas requiere acceso a la cámara. El anclaje
    de la ficha al entorno solo está disponible en dispositivos
    compatibles con ARCore.

    \item El reconocimiento automático se limita a los tipos de máquinas
    incluidos en el entrenamiento de los modelos.

    \item La generación de rutinas se basa en reglas de prescripción del
    ejercicio y en los datos del perfil del usuario.
    
\end{itemize}

